% Shared preamble for the printable cards.
%
% Each card is a standalone document that must come out as ONE A4 side. That
% constraint is the whole design: a card is consulted under pressure, and a
% card that runs to two sides is a document, which is not what a desk needs.
%
% tools/check_links.py does not police this. The build does: check-cards in the
% Makefile fails any card whose PDF has more than one page.
\documentclass[10pt,a4paper]{article}
\usepackage[T1]{fontenc}
\usepackage[utf8]{inputenc}
\IfFileExists{lmodern.sty}{\usepackage{lmodern}}{}
\usepackage[a4paper,margin=14mm]{geometry}
\usepackage{xcolor}
\usepackage{enumitem}
\usepackage{tabularx}
\usepackage{booktabs}
\usepackage{titlesec}

\definecolor{kitink}{RGB}{31,78,121}
\definecolor{kitwarn}{RGB}{150,30,30}
\definecolor{kitgrey}{RGB}{110,110,110}

\pagestyle{empty}
\setlength{\parindent}{0pt}
\setlength{\parskip}{3pt}
\setlist{itemsep=1pt, topsep=2pt, parsep=0pt, leftmargin=*}

\titleformat{\section}{\bfseries\color{kitink}}{}{0pt}{}
\titlespacing*{\section}{0pt}{7pt}{2pt}

\newcommand{\cardtitle}[2]{%
  {\LARGE\bfseries\color{kitink} #1\par}
  {\color{kitgrey}\small #2\par}
  \vspace{2pt}{\color{kitink}\rule{\linewidth}{1pt}}\par\vspace{4pt}}

\newcommand{\cardfoot}[1]{%
  \vfill{\color{kitink}\rule{\linewidth}{0.4pt}}\par\vspace{2pt}
  {\scriptsize\color{kitgrey} #1\par}}

\newcommand{\warn}[1]{\textcolor{kitwarn}{\textbf{#1}}}
% Same spelling as the book's preamble. Defined again rather than shared,
% because a card must build standalone -- somebody printing one should not
% need the book's preamble on the path.
\newcommand{\dash}{\,---\,}

\begin{document}
\cardtitle{Answering anything}{Four steps. Works on questions nobody anticipated.}

\section{1. What does it change for the person using this?}
Start at the consequence, not the mechanism. Not \textit{``a B-tree orders the
keys''} but \textit{``this scans the table, so it degrades as the table grows,
and somebody is waiting''}.\\
{\color{kitgrey}This is what makes an answer sound operational rather than
learned.}

\section{2. What does it cost to reverse?}
Sort the options by how hard they are to undo. An index: minutes. An identifier
scheme: a migration. A published API: everybody else's release cycle.\\
{\color{kitgrey}Cheap-to-reverse deserves a quick answer and a willingness to
be wrong. Expensive deserves the time. Treating them alike \dash\ either way
round \dash\ is the judgement failure.}

\section{3. What is the boring default?}
Name the obvious, well-understood option first and say why it is the default.
Then say what would make you leave it.\\
{\color{kitgrey}Starting at the interesting option has skipped the argument.}

\section{4. What would change my mind?}
The measurement, the constraint, or the fact that moves you to the other
option.\\
{\color{kitgrey}The strongest sentence in the method, and the one people leave
out. It turns a preference into a decision \dash\ and it invites them to supply
the constraint, which is often what they were waiting to do.}

\section{When you do not know}
Say where the edge is, then keep going from first principles.
\textit{``I have not used that one. What I would want to know is whether it
guarantees ordering, because the design forks there.''}\\
{\color{kitgrey}At Junior, not knowing is a gap. At Staff, knowing which
question resolves it is the skill.}

\section{The tell you are reciting}
Your answer does not change when their constraints change. If they add
``fifty users, one machine'' and the answer is the same, it was never about
their system.

\cardfoot{Staff AI Engineer Preparation Kit -- Chapter 29.}
\end{document}
